% Part: normal-modal-logic
% Chapter: axioms-systems
% Section: normal-logics

\documentclass[../../../include/open-logic-section]{subfiles}

\begin{document}

\olfileid{nml}{prf}{nor}

\olsection{నార్మల్ మోడల్ తర్కాలు}

మోడల్ \tetoken{సూత్రాల}{!!{formula}s} ప్రతి సమితినీ స్వీకృతాల
సమితి నుంచి వ్యుత్పాదించదగిన \tetoken{సూత్రాల}{!!{formula}s}
సమితిగా సులభంగా వర్ణించలేం. మోడల్ తర్కాలు చక్కగా ప్రవర్తించాలని
కోరుకుంటాం. ముందుగా, సాంప్రదాయ ప్రతిజ్ఞావాక్య తర్కంలో మనం
\tetoken{నిగమించగలిగిన}{!!{derive}} ప్రతిదీ ఇప్పటికీ
\tetoken{వ్యుత్పాదించదగినదిగా}{!!{derivable}} ఉండాలి. అయితే ఇప్పుడు
\tetoken{సూత్రాలు}{!!{formula}s},
\iftag{prvBox}{$\Box$\iftag{prvDiamond}{ మరియు~}{}}{}\iftag{prvDiamond}{$\Diamond$}{}
కూడా కలిగి ఉండవచ్చని పరిగణించాలి. అందుకోసం మోడల్ తర్కంలో
సర్వసత్య ప్రతిస్థాపన నిదర్శనాలన్నీ ఉండాలని, అది మోడస్ పోనెన్స్
కింద సంవృతమై ఉండాలని కోరుతాం.

\begin{defn}
  \emph{మోడల్ తర్కం} అంటే మోడల్ \tetoken{సూత్రాల}{!!{formula}s}
  సమితి~$\Sigma$; ఇది
  \begin{enumerate}
  \item అన్ని సర్వసత్యాలను కలిగి ఉంటుంది; మరియు
  \item ప్రతిస్థాపన కింద సంవృతమై ఉంటుంది, అంటే $!A \in \Sigma$
    అయితే, $!D_1$, \dots, $!D_n$ ఏ \tetoken{సూత్రాలైనా}{!!{formula}s},
    \[
    \SSubst{!A}{\subst{!D_1}{p_1}, \dots, \subst{!D_n}{p_n}} \in \Sigma,
    \]
    \item \emph{మోడస్ పోనెన్స్} కింద సంవృతమై ఉంటుంది, అంటే
      $!A$, $!A \lif !B \in \Sigma$ అయితే, $!B \in \Sigma$.
  \end{enumerate}
\end{defn}

మోడల్ తర్కాలకు సంబంధాత్మక అర్థవిచారాన్ని ఉపయోగించాలంటే,
అన్ని మోడల్ నమూనాల్లో చెల్లుబాటయ్యే
\tetoken{సూత్రాలన్నీ}{!!{formula}s} వాటిలో ఉండాలని కోరాలి.
\Ax{K}\iftag{prvDiamond}{ మరియు~\Dual{}}{}
యొక్క అన్ని నిదర్శనాలు \tetoken{వ్యుత్పాదించదగినవి}{!!{derivable}}
అయి, \tetoken{ఒక సూత్రం}{!!a{formula}}~$!A$
\tetoken{వ్యుత్పాదించదగినది}{!!{derivable}} అయినప్పుడల్లా
~$\Box !A$ కూడా వ్యుత్పాదించదగినదైతే ఈ షరతు నెరవేరుతుందని
తేలుతుంది. ఈ షరతులను నెరవేర్చే మోడల్ తర్కాన్ని
\emph{నార్మల్} అంటాం. (నార్మల్ కాని మోడల్ తర్కాలూ ఉన్నాయి;
అయితే సాధారణ సంబంధాత్మక నమూనాలు వాటికి సరిపోవు.)

\begin{defn}
  మోడల్ తర్కం $\Sigma$ కింది వాటిని కలిగి ఉంటే అది
  \emph{నార్మల్}:
  \begin{align*}
    \tag{\Ax{K}} & \Box(p \lif q) \lif (\Box p \lif \Box q),
    \iftag{prvDiamond}{\\
      \tag{\Dual} & \Diamond p \liff \lnot\Box\lnot p}{}
  \end{align*}
  అలాగే అది \emph{అవశ్యకీకరణ} కింద సంవృతమై ఉండాలి, అంటే
  $!A \in \Sigma$ అయితే $\Box !A \in \Sigma$.
\end{defn}

సర్వసత్య అన్వయం ఒక \emph{లోకంలో సత్యం}ను నిలిపేందుకు
తగినంత ``సూక్ష్మమైనది''. కానీ \Nec{} నియమం
\emph{ఒక నమూనాలో సత్యం}ను మాత్రమే నిలుపుతుంది (అందువల్ల
ఒక చట్రంలో లేదా చట్రాల వర్గంలో చెల్లుబాటును కూడా నిలుపుతుంది).

\begin{prop}\ollabel{prop:rk}
  ప్రతి నార్మల్ మోడల్ తర్కం \RK{} నియమం కింద సంవృతమై ఉంటుంది:
  \begin{prooftree}
    \AxiomC{$!A_1 \lif (!A_2 \lif \cdots (!A_{n-1} \lif !A_n)\cdots)$}
    \RightLabel{\RK}
    \UnaryInfC{$\Box!A_1 \lif (\Box!A_2 \lif \cdots (\Box!A_{n-1}
      \lif \Box!A_n)\cdots).$}
  \end{prooftree}
\end{prop}

\begin{proof}
  $n$పై ఆగమనం ద్వారా: $n = 1$ అయితే, ఈ నియమం \Nec{} మాత్రమే;
  ప్రతి నార్మల్ మోడల్ తర్కం \Nec{} కింద సంవృతమై ఉంటుంది.

  ఇప్పుడు ఫలితం $n-1$కు సరైందని అనుకుందాం; అది~$n$కు
  కూడా సరైందని చూపుతాం.

  అనుకుందాం:
  \begin{align*}
  & !A_1 \lif (!A_2 \lif \cdots (!A_{n-1} \lif !A_n)\cdots) \in \Sigma
  \intertext{ఆగమన పరికల్పన వల్ల మనకు}
  & \Box!A_1 \lif (\Box!A_2 \lif \cdots \Box(!A_{n-1} \lif !A_n)\cdots)
  \in \Sigma
  \intertext{$\Sigma$ నార్మల్ మోడల్ తర్కం కనుక, అది~$\Ax{K}$ యొక్క
    అన్ని నిదర్శనాలను, ముఖ్యంగా కిందిదాన్ని కలిగి ఉంటుంది:}
  & \Box(!A_{n-1} \lif !A_n) \lif (\Box!A_{n-1} \lif \Box!A_n) \in \Sigma
  \intertext{మోడస్ పోనెన్స్‌ను, తగిన సర్వసత్య ప్రతిస్థాపన
    నిదర్శనాలను ఉపయోగిస్తే మనకు}
  & \Box!A_1 \lif (\Box!A_2 \lif \cdots (\Box!A_{n-1}
  \lif \Box!A_n)\cdots) \in \Sigma.
  \end{align*}
\end{proof}

\begin{prop}\ollabel{prop:notDiamondBot}
  ప్రతి నార్మల్ మోడల్ తర్కం $\Sigma$లో~$\lnot\Diamond\lfalse$ ఉంటుంది.
\end{prop}

\begin{prob}
  \olref[nml][prf][nor]{prop:notDiamondBot}ను నిరూపించండి.
\end{prob}

\begin{prop}
  $!A_1$, \dots, $!A_n$ ఏ \tetoken{సూత్రాలైనా}{!!{formula}s}
  $!A_1$, \dots, $!A_n$ల ప్రతిస్థాపన నిదర్శనాలన్నిటినీ కలిగి ఉండే
  అతి చిన్న మోడల్
  తర్కం $\Sigma$ ఒకటి ఉంటుంది.
\end{prop}

\begin{proof}
  $!A_1$, \dots, $!A_n$ ఇచ్చినప్పుడు, $!A_1$, \dots, $!A_n$ల
  ప్రతిస్థాపన నిదర్శనాలన్నిటినీ కలిగి ఉండే అన్ని మోడల్ తర్కాల
  ప్రతిచ్ఛేదంగా $\Sigma$ను
  నిర్వచిద్దాం. అన్ని \tetoken{సూత్రాల}{!!{formula}s} సమితి
  $\Frm[L]$ కూడా అటువంటి మోడల్ తర్కమే కనుక, ప్రతిచ్ఛేదం
  తీసుకునే తర్కాల వర్గం శూన్యం కాదు. అదే వాదనలో నార్మల్
  మోడల్ తర్కాలను తీసుకుంటే, కింద ఉపయోగించే అతి చిన్న నార్మల్
  మోడల్ తర్కం లభిస్తుంది.
  \sourcecorrection{OLTENMLAXSNOR-001}{స్థిర మూల ప్రతిపాదన
    అతి చిన్న మోడల్ తర్కం గురించి చెబుతుంది; కానీ మూల నిరూపణ
    నార్మల్ మోడల్ తర్కాల ప్రతిచ్ఛేదాన్ని తీసుకొని అతి చిన్న
    నార్మల్ తర్కాన్నే చూపుతుంది. ప్రకటించిన ప్రతిపాదనకు
    అన్ని మోడల్ తర్కాల ప్రతిచ్ఛేదాన్ని వాడి, తరువాతి
    నిర్వచనానికి నార్మల్ సందర్భాన్ని వేరుగా స్పష్టం చేశాం.}
\end{proof}

\begin{defn}
$!A_1$, \dots, $!A_n$ను కలిగి ఉండే అతి చిన్న నార్మల్
మోడల్ తర్కాన్ని \emph{మోడల్ వ్యవస్థ} అంటాం; దాన్ని
$\Log{K} !A_1 \dots !A_n$తో సూచిస్తాం. అతి చిన్న
నార్మల్ మోడల్ తర్కాన్ని~\Log{K}తో సూచిస్తాం.
\end{defn}

\end{document}
